
\documentclass{article}
\usepackage{colortbl}
\usepackage{makecell}
\usepackage{multirow}
\usepackage{supertabular}

\begin{document}

\newcounter{utterance}

\centering \large Interaction Transcript for game `bbh', experiment `snarks', episode 139 with reflexion{-}Qwen3.5{-}27B.
\vspace{24pt}

{ \footnotesize  \setcounter{utterance}{1}
\setlength{\tabcolsep}{0pt}
\begin{supertabular}{c@{$\;$}|p{.15\linewidth}@{}p{.15\linewidth}p{.15\linewidth}p{.15\linewidth}p{.15\linewidth}p{.15\linewidth}}
    \# & \multicolumn{2}{c}{Player} && \multicolumn{2}{c}{Game Master} \\
    \hline

    \theutterance \stepcounter{utterance}  
    & & & \multicolumn{4}{p{0.6\linewidth}}{
        \cellcolor[rgb]{0.9,0.9,0.9}{
            \makecell[{{p{\linewidth}}}]{
                \texttt{\tiny{[P1$\langle$GM]}}
                \texttt{You have attempted this same task before and did not succeed. These are the plans you wrote after those attempts.} \\
\texttt{Trial \#1: New plan: I incorrectly identified the sarcastic statement by assuming that the option containing an illogical or irrational action (trusting a bad source) was the sarcastic one. However, in the provided examples, sarcasm was defined by the use of irony to criticize a situation, often by equating a negative reality with something absurd (like politics) or by stating a positive{-}sounding phrase that clearly means the opposite in context. In the final question, Option (A) uses the phrase "Thank you stranger on the internet" to mock the person, followed by a sincere rejection of their advice; the sarcasm lies in the polite address used to deliver a harsh critique. Option (B) is a literal statement of a foolish action, which lacks the ironic "opposite meaning" structure required by the definition. Next time, I will look for the option that uses polite or positive language to deliver a negative critique, or equates a serious issue with an absurd concept, rather than simply selecting the option that describes an illogical behavior. I will specifically analyze the tone of the opening phrase and the relationship between the premise and the conclusion to identify the ironic contrast.} \\
\\ 
\\ 
\texttt{Determine which of two sentences is sarcastic.} \\
\\ 
\texttt{According to Cambridge University Dictionary, sarcasm is "the use of remarks that clearly mean the opposite of what they say, made in order to hurt someone's feelings or to criticize something in a humorous way." Sarcastic sentences often contain satirical or ironic utterances, hyperboles, ambivalent or witty remarks.} \\
\\ 
\texttt{Q: Which statement is sarcastic?} \\
\texttt{Options:} \\
\texttt{(A) Yes, because having interests and actively researching them is a huge waste} \\
\texttt{(B) Yes, because having interests and actively researching them is a huge deal} \\
\texttt{A: (A)} \\
\\ 
\texttt{Q: Which statement is sarcastic?} \\
\texttt{Options:} \\
\texttt{(A) No one is going to disagree with you on this. Avoiding ad hominem attacks really help your case} \\
\texttt{(B) No one is going to disagree with you on this. Ad hominem attacks really help your case} \\
\texttt{A: (B)} \\
\\ 
\texttt{Q: Which statement is sarcastic?} \\
\texttt{Options:} \\
\texttt{(A) Consistency in the league's punishments? What do you think this is supposed to be, politics?} \\
\texttt{(B) Consistency in the league's punishments? What do you think this is supposed to be, moral?} \\
\texttt{A: (A)} \\
\\ 
\texttt{Q: Which statement is sarcastic?} \\
\texttt{Options:} \\
\texttt{(A) Thank you stranger on the internet, I won't trust your three word post without any sources over someone offering solid advice claiming to actually be from the country in question} \\
\texttt{(B) Thank you stranger on the internet, I will trust your three word post without any sources over someone offering solid advice claiming to actually be from the country in question} \\
\texttt{A:} \\
            }
        }
    }
     \\ \\

    \theutterance \stepcounter{utterance}  
    & \multicolumn{4}{p{0.6\linewidth}}{
        \cellcolor[rgb]{0.561,0.737,0.561}{
            \makecell[{{p{\linewidth}}}]{
                \texttt{\tiny{[P1$\rangle$GM]}}
                \texttt{Based on the definition provided and the pattern established in the previous examples, here is the analysis:} \\
\\ 
\texttt{**Analysis of the Options:**} \\
\texttt{*   **Option (A):** "Thank you stranger on the internet, I won't trust your three word post without any sources over someone offering solid advice..."} \\
\texttt{    *   The phrase "Thank you stranger on the internet" is used ironically. The speaker is not actually grateful; they are mocking the person for giving unsolicited, low{-}quality advice.} \\
\texttt{    *   The second part of the sentence ("I won't trust...") states a logical, sensible action (rejecting bad advice).} \\
\texttt{    *   The sarcasm lies in the juxtaposition of the polite "Thank you" with the harsh rejection of the advice. It fits the pattern of using a positive{-}sounding phrase to deliver a negative critique.} \\
\\ 
\texttt{*   **Option (B):** "Thank you stranger on the internet, I will trust your three word post without any sources over someone offering solid advice..."} \\
\texttt{    *   This sentence describes an illogical action (trusting a bad source over a good one).} \\
\texttt{    *   While the action is foolish, the sentence itself is a literal statement of that foolishness. It lacks the ironic "opposite meaning" structure where the words say one thing (gratitude/trust) but the context implies the opposite (mockery/rejection). In the previous examples (like the "politics" vs. "moral" or "waste" vs. "deal"), the sarcastic option was the one that used irony to highlight the absurdity, not the one that simply stated the absurdity as a fact.} \\
\\ 
\texttt{**Conclusion:**} \\
\texttt{Option (A) uses the polite address "Thank you" to mock the recipient, which aligns with the definition of sarcasm as using remarks that mean the opposite of what they say (saying "thank you" while actually criticizing) to criticize something in a humorous way.} \\
\\ 
\texttt{A: (A)} \\
            }
        }
    }
    & & \\ \\

    \theutterance \stepcounter{utterance}  
    & & & \multicolumn{2}{p{0.3\linewidth}}{
        \cellcolor[rgb]{0.9,0.9,0.9}{
            \makecell[{{p{\linewidth}}}]{
                \texttt{\tiny{[GM$|$GM]}}
                \texttt{Target: (B)} \\
            }
        }
    }
    & & \\ \\

    \theutterance \stepcounter{utterance}  
    & & & \multicolumn{2}{p{0.3\linewidth}}{
        \cellcolor[rgb]{0.9,0.9,0.9}{
            \makecell[{{p{\linewidth}}}]{
                \texttt{\tiny{[GM$|$GM]}}
                \texttt{game\_result = LOSE} \\
            }
        }
    }
    & & \\ \\

    \theutterance \stepcounter{utterance}  
    & & \multicolumn{4}{p{0.6\linewidth}}{
        \cellcolor[rgb]{0.95,0.95,0.95}{
            \makecell[{{p{\linewidth}}}]{
                \texttt{\tiny{[P1$\leftrightarrow$P]}}
                \texttt{Post{-}episode Reflexion (not sent to GM):} \\
\texttt{New plan: I previously failed because I misidentified the sarcastic option by focusing on the logical absurdity of the action rather than the ironic structure of the statement. In the provided examples, the sarcastic option consistently used a positive or polite framing to deliver a negative critique (e.g., "Thank you" while rejecting advice, or equating consistency with "politics" to imply it doesn't exist). Option (B) in the failed attempt was a literal statement of foolishness, which is not sarcasm. Option (A) used the ironic "Thank you" to mock the source while logically rejecting them, fitting the pattern of saying the opposite of what is meant to criticize. Next time, I will strictly look for the option where the surface meaning (polite/agreeing) contradicts the underlying intent (mocking/disagreeing), rather than selecting the option that simply describes an illogical or absurd scenario as a fact. I will prioritize identifying the "opposite meaning" mechanism over the content's logical validity.} \\
            }
        }
    }
    & \\ \\

\end{supertabular}
}

\end{document}
